\documentclass[conference]{IEEEtran}
\IEEEoverridecommandlockouts

\usepackage{cite}
\usepackage{amsmath,amssymb,amsfonts}
\usepackage{algorithmic}
\usepackage{graphicx}
\usepackage{textcomp}
\usepackage{xcolor}
\usepackage{booktabs}
\usepackage{url}
\usepackage{hyperref}
\usepackage{listings}
\usepackage{multirow}
\usepackage{array}

\lstset{
  basicstyle=\ttfamily\footnotesize,
  breaklines=true,
  frame=single,
  captionpos=b
}

\def\BibTeX{{\rm B\kern-.05em{\sc i\kern-.025em b}\kern-.08em
    T\kern-.1667em\lower.7ex\hbox{E}\kern-.125emX}}

\begin{document}

\title{MAEP: A Minimum Viable Agent-to-Agent Economic Protocol\\
with Decentralized Identity and Trustless Payment Settlement}

\author{
\IEEEauthorblockN{Anonymous Author(s)}
\IEEEauthorblockA{
\textit{Submitted for Review}\\
2026
}
}

\maketitle

%% ============================================================
\begin{abstract}
Existing agent communication protocols such as Google A2A and Coinbase
x402 lack two critical primitives required for fully autonomous economic
activity: verifiable decentralized agent identity and trustless
micropayment settlement with on-chain result verification.  We present
\textbf{MAEP} (Minimum Viable Agent-to-Agent Economic Protocol), a
five-stage protocol that addresses these gaps through three Ethereum
smart contracts---\texttt{AgentRegistry}, \texttt{PaymentChannel}, and
\texttt{ResultAttestation}---and a Python Agent SDK with pluggable LLM
backends.  We evaluate MAEP against x402, Google A2A, and Fetch.ai on
three dimensions: end-to-end latency overhead, on-chain Gas cost, and
trust assumption count.  Our experiments show that (1) MAEP reduces
trust assumptions to a single party (Ethereum consensus) versus 2--3
for competing protocols; (2) the total Gas cost for one complete task
delegation is 408,870 units on Base Sepolia, equating to a negligible
transaction fee at Layer-2 prices; and (3) the protocol state-machine
itself adds only microsecond-level latency above plain function calls.
An end-to-end three-agent demo---including a contested dispute
resolved by an \texttt{AuditorAgent}---completes in under five seconds.
We release the full implementation as open-source software.
\end{abstract}

\begin{IEEEkeywords}
agent-to-agent protocol, decentralized identity, smart contracts,
micropayment, trustless settlement, large language models, Web3
\end{IEEEkeywords}

%% ============================================================
\section{Introduction}
\label{sec:intro}

The emergence of large language model (LLM)-powered autonomous agents
has shifted the question of AI coordination from theoretical interest
to practical urgency.  Agents are no longer passive tools operated by
humans; they are increasingly capable of planning, tool use, and
sustained goal pursuit~\cite{park2023generative}.  A natural next step
is for agents to \emph{hire} one another: a Requester agent
decomposes a task, discovers a specialized Provider agent, delegates
work, and pays for the result---all without human intervention.

This vision requires two primitives that current protocols do not
provide together:

\begin{enumerate}
\item \textbf{Verifiable agent identity.}  How does a requester know
  that the provider it found is who it claims to be, and that its
  declared capabilities are genuine?  Centralized registries (Google
  A2A, platform-hosted agent cards) answer this by trusting the
  platform operator.  For truly autonomous, cross-organizational
  interaction, this is insufficient.

\item \textbf{Trustless payment settlement.}  How can a requester pay
  for a result without either (a) paying upfront and risking
  non-delivery, or (b) waiting for a human to authorize payment?
  Existing HTTP-based micropayment schemes (e.g., x402) rely on
  stablecoin infrastructure and server-side honesty.
\end{enumerate}

We propose \textbf{MAEP}, a Minimum Viable Agent-to-Agent Economic
Protocol that resolves both problems using public smart contracts on an
EVM-compatible blockchain.  MAEP is deliberately minimal: it does the
least possible on-chain work necessary to remove the need for a trusted
intermediary, while delegating expensive computation (LLM inference,
task matching) to off-chain agents.

The main contributions of this paper are:

\begin{itemize}
\item A formal five-stage protocol specification (Register, Discover,
  Delegate, Execute, Settle) with defined state transitions, message
  formats, and security invariants (\S\ref{sec:protocol}).

\item A reference implementation comprising three audited Solidity
  contracts and a Python SDK supporting six LLM providers
  (\S\ref{sec:implementation}).

\item A comparative evaluation on latency, Gas cost, and trust
  assumption count, demonstrating that MAEP achieves the lowest trust
  footprint among surveyed protocols (\S\ref{sec:evaluation}).
\end{itemize}

%% ============================================================
\section{Background}
\label{sec:background}

This section introduces the foundational concepts required to
understand MAEP.  Readers familiar with blockchain and A2A protocols
may skim to \S\ref{sec:related}.

\subsection{Blockchain and Smart Contracts}

A \emph{blockchain} is a distributed ledger replicated across thousands
of nodes.  Transactions are grouped into blocks, and each block
cryptographically references its predecessor, forming a tamper-evident
chain.  The key property relevant here is \emph{trustless execution}:
any computation defined in a smart contract executes identically on
every node, and the result is final once confirmed.

\emph{Smart contracts} are programs stored on-chain.  They have a
persistent address, can hold cryptocurrency balances, and expose
functions callable by any external account.  Crucially, their source
code is public and their execution is deterministic---no operator can
alter behavior after deployment.

The \emph{Ethereum Virtual Machine} (EVM) is the standard runtime
environment for smart contracts.  An instruction unit called
\emph{Gas} prices each EVM operation; users pay Gas fees in the
network's native currency (ETH) to execute transactions.  Gas prices
fluctuate with network demand, but Gas \emph{units} consumed by a
given computation are deterministic---the same code path always costs
the same number of Gas units regardless of when or where it is
executed.

\textbf{Layer 2 (L2) networks} (such as Base, Arbitrum, and Optimism)
run the same EVM but batch transactions and post proofs to Ethereum
mainnet, reducing fees by two to three orders of magnitude.
MAEP targets \textbf{Base}, a Coinbase-operated Optimistic Rollup,
because its tooling is mature and its gas prices are typically below
0.01 gwei---making sub-cent transaction fees realistic.

\subsection{Cryptographic Hash Functions}

A \emph{hash function} maps arbitrary-length data to a fixed-length
digest.  The \texttt{keccak256} function used in Ethereum is a
one-way function: it is computationally infeasible to recover the
input from the digest, and any change to the input changes the digest
completely.  MAEP uses \texttt{keccak256(result\_data)} as a
\emph{commitment}: the provider posts only the hash on-chain (cheap),
and later reveals the full result off-chain.  The requester verifies
that \texttt{keccak256(received\_data) == on\_chain\_hash}.

\subsection{Escrow and Payment Channels}

An \emph{escrow} is a neutral holding mechanism: funds are locked by
a third party and released only when predefined conditions are met.
Traditional escrow requires a trusted third party (a bank, a lawyer).
Smart-contract escrow replaces the trusted party with code: funds are
locked in the contract's balance and released when the verification
logic returns true.

MAEP's \texttt{PaymentChannel} contract implements a simple escrow:
(1) the requester calls \texttt{lock(taskId, provider, deadline)} with
attached ETH; (2) if the provider submits a result and the requester
confirms, \texttt{release(taskId)} transfers funds to the provider;
(3) if the deadline passes without settlement, \texttt{refund(taskId)}
returns funds to the requester.  No human intervention is required at
any step.

\subsection{Decentralized Identity (DID)}

\emph{Decentralized Identifiers} (DIDs) are globally unique identifiers
controlled by their subject rather than by a registrar.  The W3C DID
specification~\cite{w3c_did} defines a data model; different DID
methods specify how identifiers are created and resolved on different
substrates (IPFS, Bitcoin, Ethereum, etc.).

MAEP implements a lightweight on-chain DID registry.  Each agent
registers a JSON-encoded capability document
(e.g., \texttt{\{"task\_types":["data\_analysis"]\}}) alongside a
minimum stake in ETH.  The stake creates a mild Sybil-resistance
mechanism: an attacker registering $N$ fake identities must lock
$N \times \text{MIN\_STAKE}$ ETH.  Reputation scores are updated
on-chain after each settlement, providing a persistent, tamper-evident
track record.

\subsection{Agent-to-Agent (A2A) Communication}

An \emph{A2A protocol} defines how autonomous agents discover each
other, negotiate tasks, and exchange results.  The essential components
are:

\begin{itemize}
\item \textbf{Discovery}: how agents find peers with needed capabilities.
\item \textbf{Task specification}: a shared format for describing what
  is needed and at what price.
\item \textbf{Result delivery}: how the provider returns its output.
\item \textbf{Settlement}: how payment flows after task completion.
\end{itemize}

Current proposals (Google A2A, x402, Fetch.ai) provide subsets of
these components.  MAEP aims to be the minimal protocol that covers
all four with trust-minimized guarantees.

\subsection{Large Language Model Agents}

An \emph{LLM agent} wraps a language model with a control loop that
allows it to use tools, maintain state, and pursue multi-step goals.
In MAEP's reference implementation, agents are Python objects that
hold an \texttt{LLMClient} (which abstracts over different providers)
and a \texttt{MAEPSession} (which tracks protocol state).  When a
provider agent receives a task, it constructs a prompt, calls its LLM
backend, and returns the generated text as the result.  The SDK is
deliberately LLM-agnostic: the same agent code works with OpenAI GPT,
Anthropic Claude, local Ollama models, or any OpenAI-compatible
endpoint.

%% ============================================================
\section{Related Work}
\label{sec:related}

\subsection{Google Agent2Agent (A2A) Protocol}

\subsubsection{Overview and Motivation}

Google released the Agent2Agent (A2A) protocol in April 2025 as an
open standard for interoperability between AI agents built on different
platforms~\cite{google_a2a}.  The driving observation is that
enterprise AI deployments increasingly involve agents from multiple
vendors (e.g., a Salesforce CRM agent delegating to a SAP logistics
agent), and without a shared communication standard each pair of
vendors must negotiate a custom integration.  A2A defines a common
JSON-over-HTTP wire format so that any compliant agent can talk to any
other without pre-arranged bilateral agreements at the
\emph{communication} level.

\subsubsection{Core Mechanism: Agent Cards and Task Lifecycle}

The central abstraction in A2A is the \textbf{Agent Card}---a
JSON document served at a well-known HTTPS endpoint
(e.g., \texttt{https://agent.example.com/.well-known/agent.json}).
An Agent Card declares:

\begin{itemize}
  \item \textbf{Identity fields}: agent name, description, version,
    and the URL of the agent's A2A endpoint.
  \item \textbf{Capability list}: a set of \emph{skills}, each with a
    natural-language description, input/output schema, and optional
    example messages.
  \item \textbf{Authentication requirements}: which OAuth 2.0 scopes
    or API keys the caller must present.
  \item \textbf{Supported modalities}: whether the agent can handle
    text, file attachments, structured data, audio, or video.
\end{itemize}

A \textbf{Requester} agent (called the \emph{client} in A2A
terminology) discovers a Provider agent (the \emph{remote agent}) by
fetching its Agent Card, inspects the declared skills, and initiates a
\textbf{Task}.  The task lifecycle follows four states:

\begin{enumerate}
  \item \texttt{submitted} -- the client sends a \texttt{tasks/send}
    JSON-RPC call containing the task input and a unique task ID.
  \item \texttt{working} -- the remote agent is processing; it may
    stream intermediate \emph{artifacts} (partial results) back to the
    client via Server-Sent Events (SSE) or WebSockets.
  \item \texttt{input-required} -- the remote agent needs clarification
    and pauses to request additional user or agent input (supporting
    multi-turn interaction).
  \item \texttt{completed} / \texttt{failed} / \texttt{canceled} --
    terminal states returned with final artifacts or error codes.
\end{enumerate}

All messages are JSON objects transported over HTTPS, making A2A easy
to implement in any language and compatible with existing API
infrastructure.

\subsubsection{Identity and Trust Model}

\textbf{How A2A establishes identity.}
In A2A, an agent's identity is anchored to its \emph{HTTPS domain}.
When Agent~B serves its Agent Card at
\texttt{https://b.example.com/.well-known/agent.json}, Agent~A trusts
that the card genuinely represents the operator of
\texttt{b.example.com} because HTTPS certificate validation
guarantees the domain owner controls the TLS private key.  This is the
same trust model as ordinary Web PKI: the browser trusts a certificate
because a Certificate Authority (CA) has vouched for the domain.

\textbf{Limitations of domain-anchored identity.}
This approach has three concrete weaknesses relevant to autonomous
agent economies:

\begin{enumerate}
  \item \textbf{Platform operator trust.}  For agents hosted on a
    managed platform (e.g., Google Agentspace, Microsoft Copilot
    Studio), the Agent Card is served by the \emph{platform}, not by
    the agent itself.  If the platform operator modifies the declared
    capabilities, censors an agent, or revokes access, the requester
    has no cryptographic means to detect the change.  The requester
    trusts the platform operator by necessity.

  \item \textbf{Capability claims are unverified.}  An Agent Card
    declares skills in natural language.  There is no mechanism in the
    A2A protocol to \emph{attest} that a provider has successfully
    completed tasks of the declared type in the past, or to penalize
    it for false claims.  A malicious provider can register arbitrary
    capabilities with no cost and no accountability.

  \item \textbf{No cross-organizational reputation.}  A2A provides no
    persistent reputation or history mechanism.  If an agent from
    Company~X fails to deliver on a task for Company~Y, Company~Y has
    no standardized way to record this failure in a location visible to
    future requesters from Company~Z.  Each organization must maintain
    its own private allow/deny list.
\end{enumerate}

\textbf{Authentication (not the same as identity).}
A2A supports OAuth 2.0, API keys, and other standard HTTP
authentication schemes.  This authenticates \emph{who is calling} at
the transport level, but does not address the deeper question of
whether the callee's declared capabilities are trustworthy.
Authentication proves ``this request comes from account X''; it does
not prove ``account Y (the callee) has the skills it claims and will honor its
commitments.''

\subsubsection{Economic Primitives: Deliberately Absent}

A2A is explicitly scoped to \emph{communication}, not economics.  The
protocol specification states that payment is out of scope and should
be handled by ``existing payment infrastructure.''  Concretely:

\begin{itemize}
  \item There is \textbf{no escrow}: a requester who pre-pays has no
    protocol-level recourse if the provider returns garbage.
  \item There is \textbf{no result verification}: the protocol
    delivers artifacts but does not commit to their hash or provide
    any proof of correct computation.
  \item There is \textbf{no dispute resolution}: the specification does
    not define what happens when the requester rejects a result.
\end{itemize}

This is a reasonable scope decision for A2A's primary use case
(enterprise integration within trusted organizational boundaries), but
it means A2A is insufficient as a substrate for open, permissionless
agent markets where requesters and providers may have no prior
relationship and no shared institutional accountability.

\subsubsection{Comparison with MAEP}

MAEP treats A2A as a \emph{complementary} protocol rather than a
competitor.  A2A's message format and task lifecycle could serve as
the off-chain communication layer for MAEP's DISCOVER and EXECUTE
stages---agents would find each other and exchange task details using
A2A messages, while using MAEP's smart contracts for binding identity
registration, payment escrow, and result attestation.  In this view,
MAEP adds the \emph{economic trust layer} that A2A deliberately omits.

\subsection{Coinbase x402}

x402~\cite{x402} extends HTTP with a \texttt{402 Payment Required}
response and a standardized payment header.  The client prepares an
on-chain USDC transfer; the server verifies the transaction and
proceeds.  x402 solves the \emph{pay-per-request} problem elegantly
for server-side APIs, but it does not provide agent identity, result
verification, or dispute resolution.  Payment flows from client to
server without any escrow; if the server misbehaves, the client has no
recourse within the protocol.

\subsection{Fetch.ai / Agentverse}

Fetch.ai~\cite{fetchai} provides a full ecosystem for autonomous
economic agents (AEAs) built on its own Cosmos-based blockchain.
Agents use the FET token for payment, and the network handles identity
and discovery.  Fetch.ai is more feature-complete than MAEP for
in-ecosystem agents, but it requires trust in the Fetch.ai validator
set and governance, which is a permissioned subset of the broader
Ethereum trust base.  MAEP runs on Ethereum/Base, which has a larger,
more decentralized validator set.

\subsection{Autonolas / Olas}

Autonolas~\cite{autonolas} defines on-chain service registries and
multi-agent service coordination, with staking and slashing for
misbehaving agents.  Olas is a production system with real staking
economics.  MAEP is complementary: where Olas focuses on complex
service orchestration with DAOs, MAEP targets the simpler bilateral
task delegation primitive in a minimally invasive way.

\subsection{Recent Academic Work}

Closest to our work is a concurrent preprint by Andres et
al.~\cite{andres2025}, which proposes extending A2A with blockchain
identity and payments.  Their design is conceptually aligned with
ours; MAEP differs in providing (a) a complete formal protocol
specification, (b) a working open-source implementation with
quantitative experiments, and (c) an explicit trust-minimization
analysis.

%% ============================================================
\section{MAEP Protocol Specification}
\label{sec:protocol}

\subsection{Threat Model and Design Goals}

We consider two rational, potentially self-interested parties:
\textbf{Requester} and \textbf{Provider}.  Neither is assumed to be
honest, but both are assumed to be economically rational (they will
not take actions that cost them more than they gain).

\textbf{Threat model:}
\begin{itemize}
\item Provider may claim false capabilities at registration.
\item Provider may return garbage results after receiving payment.
\item Requester may refuse to acknowledge a valid result to avoid payment.
\item Either party may abandon the protocol mid-execution.
\end{itemize}

\textbf{Design goals:}
\begin{enumerate}
\item \textbf{G1 -- Payment safety}: the requester cannot lose funds
  without receiving a committed result hash.
\item \textbf{G2 -- Liveness}: if both parties act honestly, the
  protocol always terminates in SETTLED.
\item \textbf{G3 -- Sybil resistance}: creating fake identities has a
  non-zero cost.
\item \textbf{G4 -- Minimal trust}: the only trusted party is the
  EVM's execution guarantee.
\end{enumerate}

\subsection{Five-Stage Protocol Flow}

Table~\ref{tab:stages} summarizes the five stages.  Figure~\ref{fig:flow}
illustrates the message sequence.

\begin{table}[h]
\caption{MAEP Protocol Stages}
\label{tab:stages}
\centering
\begin{tabular}{llll}
\toprule
\textbf{Stage} & \textbf{Actor} & \textbf{On-chain?} & \textbf{Key Action} \\
\midrule
1. REGISTER  & Both     & Yes & Submit DID + capabilities + stake \\
2. DISCOVER  & Requester & No  & Match by task type + budget \\
3. DELEGATE  & Requester & Yes & Lock ETH in escrow; send task spec \\
4. EXECUTE   & Provider  & Yes & Submit \texttt{keccak256(result)} \\
5. SETTLE    & Both      & Yes & Verify; release payment; update rep. \\
\bottomrule
\end{tabular}
\end{table}

\textbf{Stage 1 -- REGISTER.}
An agent calls \texttt{AgentRegistry.register(capabilityJSON)} with a
minimum stake of 0.0001~ETH attached.  The contract stores the agent's
address, capability document, reputation score (initialized to 100),
and stake.  The stake is returned on deregistration if reputation
remains above a threshold.

\textbf{Stage 2 -- DISCOVER.}
The requester queries the off-chain Python SDK's \texttt{discover()}
method, which reads the on-chain registry and filters agents whose
declared \texttt{task\_types} include the required type and whose
reputation exceeds a configurable minimum.  Discovery is off-chain to
avoid Gas overhead for what may be a frequently repeated read
operation.  The on-chain registry remains the authoritative source of
truth.

\textbf{Stage 3 -- DELEGATE.}
The requester calls
\texttt{PaymentChannel.lock(taskId, providerAddress, deadline)}
with the agreed payment attached as ETH.  The contract records:
\[\texttt{locks[taskId]} = (\text{requester}, \text{provider}, \text{amount}, \text{deadline}, \texttt{LOCKED})\]
The requester then sends the full task specification to the provider
over any off-chain channel (HTTP, WebSocket, etc.).  The \texttt{taskId}
serves as the binding link between the off-chain task and the on-chain
escrow.

\textbf{Stage 4 -- EXECUTE.}
The provider processes the task using its LLM backend and calls
\texttt{ResultAttestation.attest(taskId, keccak256(resultData))}.
This commits the result on-chain before revealing the full data,
preventing the provider from changing the result after seeing the
requester's reaction (\emph{commit-reveal pattern}).  The provider
then delivers \texttt{resultData} to the requester off-chain.

\textbf{Stage 5 -- SETTLE.}
The requester verifies that \texttt{keccak256(received)} matches the
on-chain attestation.  If accepted, it calls
\texttt{PaymentChannel.release(taskId)}, which transfers ETH to the
provider and increments the provider's reputation.  If rejected, the
session enters DISPUTED state and escalates to an \texttt{AuditorAgent},
which independently evaluates both the task specification and the
result and issues a ruling.  If no settlement occurs before the
deadline, either party can call \texttt{PaymentChannel.refund(taskId)}
to recover the locked funds.

\subsection{Security Analysis}

\textbf{G1 (Payment safety).}
Funds are transferred only in \texttt{release()} or \texttt{refund()}.
The contract follows the Checks-Effects-Interactions (CEI)
pattern~\cite{consensys_patterns}: state is updated before any
external call, preventing reentrancy attacks.

\textbf{G2 (Liveness).}
If both parties act honestly, the provider submits a valid result, the
requester's verification passes, and \texttt{release()} is called.
Two failure modes are covered by the protocol:

\emph{Unresponsive provider.}  If the provider never submits a result,
the deadline mechanism protects the requester: after the deadline
expires, the requester calls \texttt{refund(taskId)} to recover the
locked ETH.  No human intervention is required; the contract enforces
the timeout autonomously.

\emph{Unresponsive requester.}  If the provider submits a valid result
hash but the requester neither calls \texttt{release()} nor
\texttt{dispute()}, the funds remain locked until the deadline passes,
at which point the requester may reclaim them via \texttt{refund()}.
This is a known asymmetry in MAEP's current design: a requester can
effectively withhold payment without triggering a dispute, so long as
it accepts the deadline penalty.  Mitigations include requiring the
requester to explicitly reject (triggering mandatory arbitration) or
implementing an auto-release after a response deadline---both are left
as future work (see \S\ref{sec:discussion}).

\textbf{G3 (Sybil resistance).}
Each registration requires a minimum stake.  An adversary creating
$N$ identities locks $N \times 0.0001$~ETH, providing economic
friction without prohibitive cost for legitimate agents.

\textbf{G4 (Minimal trust).}
Every conditional fund transfer is governed by contract logic, not by
any operator.  The \texttt{AuditorAgent} is the only subjective
component; its ruling could be replaced by a DAO vote or a
cryptographic oracle without changing the protocol's on-chain
structure.

%% ============================================================
\section{Implementation}
\label{sec:implementation}

\subsection{Smart Contracts}

The three contracts are implemented in Solidity 0.8.24 and compiled
with the Hardhat 3 development framework.  The full source is
available in the open-source repository.

\textbf{AgentRegistry.sol} (key interface):
\begin{lstlisting}[caption={AgentRegistry interface}]
function register(string calldata capabilities)
    external payable;

function updateReputation(address agent, int8 delta)
    external;   // callable by PaymentChannel only

function deregister() external;
\end{lstlisting}

Reputation is stored as a \texttt{uint16} bounded to [0, 10000].
The \texttt{updateReputation} function is access-controlled to the
\texttt{PaymentChannel} address set at deployment, preventing arbitrary
manipulation.

\textbf{PaymentChannel.sol} (key interface):
\begin{lstlisting}[caption={PaymentChannel interface}]
function lock(bytes32 taskId,
              address provider,
              uint256 deadline)
    external payable;

function release(bytes32 taskId) external;

function refund(bytes32 taskId) external;
\end{lstlisting}

ETH transfers use the low-level \texttt{.call\{value\}("")} pattern
rather than \texttt{.transfer()} to avoid gas stipend issues with
smart contract recipients.

\textbf{ResultAttestation.sol} (key interface):
\begin{lstlisting}[caption={ResultAttestation interface}]
function attest(bytes32 taskId, bytes32 resultHash)
    external;

function dispute(bytes32 taskId) external;

function getAttestation(bytes32 taskId)
    external view
    returns (bytes32 hash, address attester,
             uint256 timestamp, bool disputed);
\end{lstlisting}

\subsection{Python Agent SDK}

The SDK is structured around four modules:

\begin{itemize}
\item \texttt{config.py}: reads \texttt{LLM\_PROVIDER},
  \texttt{LLM\_API\_KEY}, \texttt{LLM\_MODEL}, and
  \texttt{LLM\_BASE\_URL} from a \texttt{.env} file.

\item \texttt{llm\_client.py}: a unified wrapper supporting OpenAI,
  Anthropic, Zhipu, MiniMax, Ollama, and any OpenAI-compatible
  endpoint.  The \texttt{complete(prompt)} method transparently
  handles provider-specific authentication and response parsing,
  including fallback from empty \texttt{content} to \texttt{reasoning}
  fields for chain-of-thought models.

\item \texttt{protocol.py}: the \texttt{MAEPSession} state machine and
  \texttt{TaskSpec} / \texttt{TaskResult} dataclasses.

\item \texttt{requester.py}, \texttt{provider.py},
  \texttt{auditor.py}: agent roles that compose an
  \texttt{LLMClient} with protocol operations.
\end{itemize}

\subsection{Deployment}

Contracts are deployed using a Hardhat deployment script with
\texttt{ethers.js}.  Deployed addresses on Base Sepolia (chain ID
84532) are:

\begin{itemize}
\item \texttt{AgentRegistry}: \texttt{0xc4146C22...d7445}
\item \texttt{PaymentChannel}: \texttt{0x4fA1a014...7794}
\item \texttt{ResultAttestation}: \texttt{0xC915A7e9...f461}
\end{itemize}

%% ============================================================
\section{Evaluation}
\label{sec:evaluation}

We evaluate MAEP along four axes, corresponding to the four
experiments described below.  All experiments were run on a Windows 11
machine (Intel i7, 16 GB RAM) with Hardhat 3 for local measurements
and the public Base Sepolia RPC endpoint for on-chain measurements.

\subsection{Experiment 1: Latency Overhead}

\textbf{Goal.}
Quantify the latency of MAEP's protocol operations at two levels:
(1) off-chain state-machine overhead versus equivalent bare Python
operations; and (2) on-chain confirmation latency per stage on a
local Hardhat network.

\textbf{Setup.}
We use a two-table design to separate the two fundamentally different
latency sources.

\emph{Table~\ref{tab:latency_offchain} -- Off-chain overhead.}
We ran $n=200$ iterations of each stage in two configurations.
\textbf{MAEP}: each stage is timed in isolation on a pre-built session
(prior stages completed but not re-timed), so measurements are not
cumulative.  \textbf{Baseline}: the equivalent bare Python operations
that a plain HTTP A2A client would perform at the same stage---
specifically, dict construction for DELEGATE, a \texttt{hashlib.sha256}
call plus dict assignment for EXECUTE (matching MAEP's hash
commitment step), and a boolean branch plus dict update for SETTLE.
The baseline thus reflects the irreducible work any protocol must do at
each stage, making the comparison apples-to-apples rather than
MAEP-vs.-nothing.  All measurements use \texttt{time.perf\_counter\_ns()}
and are local in-process to isolate from network and blockchain latency.

\emph{Table~\ref{tab:latency_onchain} -- On-chain confirmation latency.}
We ran $n=5$ full MAEP task lifecycles against a local Hardhat node,
timing each stage from \texttt{send\_raw\_transaction()} to
\texttt{wait\_for\_transaction\_receipt()}.  This captures real EVM
execution plus JSON-RPC round-trip overhead, and represents the minimum
achievable on-chain latency (no real network propagation).

\textbf{Results.}
Tables~\ref{tab:latency_offchain} and~\ref{tab:latency_onchain} report
the results.

\begin{table}[h]
\caption{Off-chain protocol overhead per stage ($n=200$, $\mu$s).
         Each stage is timed in isolation on a pre-built session.
         Baseline = equivalent bare Python operations (dict, sha256, branch).}
\label{tab:latency_offchain}
\centering
\begin{tabular}{lrrrr}
\toprule
\textbf{Stage} & \textbf{MAEP mean} & \textbf{Base mean} & \textbf{Overhead} & \textbf{MAEP $\sigma$} \\
\midrule
DELEGATE & 1.12 & 0.23 & 4.8$\times$ & 0.47 \\
EXECUTE  & 3.32 & 1.08 & 3.1$\times$ & 2.91 \\
SETTLE   & 1.47 & 0.22 & 6.8$\times$ & 0.27 \\
\bottomrule
\end{tabular}
\end{table}

\begin{table}[h]
\caption{On-chain confirmation latency on Hardhat local network
         ($n=5$, ms, mean $\pm$ std).}
\label{tab:latency_onchain}
\centering
\begin{tabular}{lrrr}
\toprule
\textbf{Stage} & \textbf{Mean (ms)} & \textbf{Std (ms)} & \textbf{Range (ms)} \\
\midrule
REGISTER & 11.2 & 9.5  & 5.1--28.0 \\
DELEGATE &  7.8 & 2.5  & 6.0--12.2 \\
EXECUTE  &  5.5 & 0.5  & 4.6--6.1  \\
SETTLE   &  6.0 & 0.4  & 5.5--6.4  \\
\midrule
\textbf{TOTAL} & \textbf{30.5} & -- & -- \\
\bottomrule
\end{tabular}
\end{table}

\textbf{Analysis.}
\emph{Off-chain overhead (Table~\ref{tab:latency_offchain}).}
The MAEP state machine adds 1--3.3~$\mu$s per stage above the
equivalent bare operations.  The EXECUTE stage is the most expensive
because it both computes a hash and performs a dataclass instantiation;
its high standard deviation ($\sigma = 2.91~\mu$s) reflects CPU cache
effects across 200 runs.  The SETTLE stage shows the highest relative
overhead (6.8$\times$) because its baseline---a boolean branch---is
extremely cheap; in absolute terms the overhead is only 1.25~$\mu$s.
All values are sub-microsecond to low-microsecond, confirming that the
protocol state machine is not a performance bottleneck.

\emph{On-chain confirmation latency (Table~\ref{tab:latency_onchain}).}
On Hardhat local, the four on-chain stages complete in 30.5~ms total on
average.  The REGISTER stage has the highest variance (std~9.5~ms)
because it involves a larger storage write (JSON capability string),
causing occasional JIT compilation effects in the EVM interpreter.
On a real network such as Base Sepolia, each stage requires waiting for
a block (typically 2 seconds), making the total per-task on-chain
latency approximately 8 seconds---four to five orders of magnitude
larger than the off-chain state-machine overhead measured in
Table~\ref{tab:latency_offchain}.  This confirms that MAEP's design
decision to perform DISCOVER off-chain is well-motivated: moving even
one additional read operation on-chain would add seconds of latency
with no security benefit.

\subsection{Experiment 2: Gas Cost Analysis}

\textbf{Goal.}
Measure on-chain Gas consumption per stage to assess economic
viability and guide future optimization.

\textbf{Setup.}
We ran the full MAEP five-stage flow on both a local Hardhat network
(deterministic, zero-latency) and the Base Sepolia testnet (chain ID
84532).  Because EVM Gas units are fully deterministic for a given
bytecode and call data, local and testnet values are expected to
match.  We report both for validation.

\textbf{Results.}
Table~\ref{tab:gas} reports Gas units per stage.

\begin{table}[h]
\caption{Gas consumption per stage: local Hardhat vs.\ Base Sepolia}
\label{tab:gas}
\centering
\begin{tabular}{lrrl}
\toprule
\textbf{Stage} & \textbf{Hardhat} & \textbf{Base Sepolia} & \textbf{Match} \\
\midrule
REGISTER & 138,625 & 138,625 & Exact \\
DELEGATE & 116,489 & 116,501 & $\Delta$+12 \\
EXECUTE  & 93,590  & 93,602  & $\Delta$+12 \\
SETTLE   & 62,642  & 60,142  & $\Delta$-2,500 \\
\midrule
\textbf{TOTAL} & \textbf{411,346} & \textbf{408,870} & $\approx$0.6\% diff \\
\bottomrule
\end{tabular}
\end{table}

Table~\ref{tab:deploy_gas} reports one-time contract deployment costs.

\begin{table}[h]
\caption{Contract deployment Gas (one-time)}
\label{tab:deploy_gas}
\centering
\begin{tabular}{lr}
\toprule
\textbf{Contract} & \textbf{Gas Units} \\
\midrule
AgentRegistry    & 1,234,274 \\
PaymentChannel   & 1,161,550 \\
ResultAttestation & 516,032 \\
\midrule
Total deployment & 2,911,856 \\
\bottomrule
\end{tabular}
\end{table}

\textbf{Analysis.}
Gas consumption decreases monotonically from REGISTER to SETTLE,
reflecting the storage-write-heavy nature of early stages and the
storage-read-plus-transfer nature of the final stage.  The small
discrepancies between local and testnet values ($<$1\%) arise from
minor differences in block context (e.g., BASEFEE) that affect a
small number of opcodes; they validate the equivalence of local
development as a proxy for testnet behavior.

At the Base Sepolia gas price observed during our experiment
(0.006 gwei), the full task lifecycle costs:
\[408{,}870 \times 0.006 \times 10^{-9} \text{ ETH} \approx 2.45 \times 10^{-6} \text{ ETH}\]
At an ETH price of \$3{,}000 USD, this is approximately \textbf{\$0.0074
per task delegation}---well within the range of commercial viability
for agent services priced at cents to dollars.

The per-stage Gas breakdown also reveals a useful design validation:
the REGISTER stage is the most expensive because it writes a variable-length
JSON string on-chain, while SETTLE is cheapest because it mainly
reads existing state and transfers ETH.  This aligns with our design
goal of minimizing per-task overhead by amortizing registration costs
across many task executions.

\subsection{Experiment 3: Trust Assumption Analysis}

\textbf{Goal.}
Formally enumerate and compare the trust assumptions required by MAEP,
x402, Google A2A, and Fetch.ai.

\textbf{Method.}
For each protocol we enumerate trust requirements along four
dimensions: agent identity, payment, result verification, and dispute
resolution.  A ``trust assumption'' is any entity whose honesty or
availability must be assumed for the protocol's security properties to
hold.

\textbf{Results.}
Table~\ref{tab:trust} summarizes the comparison.

\begin{table}[h]
\caption{Trust assumption comparison across A2A protocols}
\label{tab:trust}
\centering
\renewcommand{\arraystretch}{1.15}
\begin{tabular}{p{1.5cm}p{1.5cm}p{1.5cm}p{1.5cm}p{1.5cm}}
\toprule
\textbf{Dimension} & \textbf{x402} & \textbf{Google A2A} & \textbf{Fetch.ai} & \textbf{MAEP} \\
\midrule
Agent Identity
  & HTTP server
  & Google platform
  & Validator set
  & Ethereum contract (public) \\
\addlinespace
Payment
  & Stablecoin issuer + server
  & None
  & FET + validators
  & Ethereum contract \\
\addlinespace
Result Verification
  & None
  & None
  & Fetch.ai network
  & Ethereum contract \\
\addlinespace
Dispute Resolution
  & None
  & Platform operator
  & Fetch.ai governance
  & AuditorAgent (configurable) \\
\midrule
\textbf{Trust Count} & \textbf{3} & \textbf{2} & \textbf{2} & \textbf{1} \\
\bottomrule
\end{tabular}
\end{table}

\textbf{Analysis.}
MAEP achieves the lowest trust assumption count (1) among the surveyed
protocols.  The single remaining assumption---trust in Ethereum
consensus---is shared with the entirety of the Ethereum DeFi
ecosystem, which has processed trillions of dollars in value.  It is a
public, permissionless trust base with over 900,000 active validators
as of 2025.

x402 requires three distinct trusted parties: the HTTP server (for
honest accounting), the stablecoin issuer (for USDC redemption
guarantees), and implicitly the payment processor routing the
transaction.  A compromise of any one breaks payment safety.

Google A2A's platform trust assumption is qualitatively different from
blockchain trust: it relies on Google's organizational integrity and
continued operation, which are subject to business and regulatory
risk.

Fetch.ai's validator set trust is formally similar to Ethereum's, but
the set is smaller (permissioned), making it more susceptible to
coordinated attacks.  Furthermore, Fetch.ai's governance assumption
adds a social-consensus dependency that MAEP avoids.

The \texttt{AuditorAgent} in MAEP's dispute path is acknowledged as a
trust assumption, but it is architecturally pluggable---it could be
replaced by a multi-sig DAO, a cryptographic proof system, or
eliminated for protocols where disputes trigger automatic refund.

\subsection{Experiment 4: End-to-End Scenario Demonstration}

\textbf{Goal.}
Demonstrate that MAEP's SDK enables a real three-agent system to
execute both a cooperative and a contested task delegation end-to-end.

\textbf{Setup.}
Three agents are instantiated: \texttt{RequesterAgent} (issues task),
\texttt{ProviderAgent} (executes with LLM backend), and
\texttt{AuditorAgent} (resolves disputes).  All agents use a locally
running \texttt{gpt-oss:20b-cloud} model via the Ollama interface.
Two scenarios are run:

\begin{itemize}
\item \textbf{Scenario A (Happy Path)}: The requester delegates a
  data analysis task (``Analyze: 100 users, avg age 32, 60\% male.
  Provide 3 key insights.'').  The provider generates insights, posts
  the result hash on-chain, and the requester verifies and accepts.

\item \textbf{Scenario B (Dispute Path)}: The requester delegates
  analysis of an empty dataset.  The provider returns a technically
  valid response (''dataset is empty---no analysis possible'').  The
  requester rejects the result.  The \texttt{AuditorAgent} reviews
  both the task specification and the result and issues a ruling.
\end{itemize}

\textbf{Results.}

\begin{table}[h]
\caption{End-to-end scenario outcomes}
\label{tab:scenario}
\centering
\begin{tabular}{llrr}
\toprule
\textbf{Scenario} & \textbf{Final Stage} & \textbf{Duration (s)} & \textbf{Ruling} \\
\midrule
A: Happy Path & SETTLED  & 1.84 & -- \\
B: Dispute    & DISPUTED & 4.20 & PROVIDER \\
\bottomrule
\end{tabular}
\end{table}

\textbf{Analysis.}
Scenario A completes in 1.84 seconds end-to-end, the majority of which
is LLM inference time.  Protocol overhead (state transitions, hash
computation) is unmeasurably small relative to inference latency,
consistent with Experiment 1.

In Scenario B, the \texttt{AuditorAgent} rules in favor of the
\textbf{Provider} (``PROVIDER'' ruling), correctly identifying that
the provider's response is technically valid given the empty dataset,
and that the requester's rejection is unreasonable.  This demonstrates
that the dispute-resolution layer can exercise independent judgment
rather than simply siding with the payer.  The 4.20-second duration
reflects two sequential LLM calls (provider execution + auditor
evaluation).

Together, the two scenarios confirm that MAEP's SDK correctly
implements the state machine transitions DELEGATED $\to$ EXECUTED
$\to$ SETTLED and DELEGATED $\to$ EXECUTED $\to$ DISPUTED, and that
the \texttt{AuditorAgent}'s LLM-driven arbitration produces sensible
results on a qualitative task.

%% ============================================================
\section{Discussion}
\label{sec:discussion}

\subsection{Limitations}

\textbf{Off-chain result delivery.}
MAEP commits a result \emph{hash} on-chain but delivers the full
result off-chain.  If the provider sends a hash that does not match
the off-chain data, the requester can detect the mismatch and dispute,
but cannot recover the correct result on-chain.  A future extension
could use IPFS content-addressed storage so that
\texttt{ipfsCID(result)} is simultaneously the hash and the retrieval
address.

\textbf{Integrity vs.\ correctness.}
MAEP's \texttt{ResultAttestation} contract verifies \emph{integrity}:
it guarantees that the result the requester received is exactly what
the provider committed on-chain, with no post-hoc tampering.  It does
\emph{not} verify \emph{correctness}: a provider can submit a
cryptographically valid hash whose pre-image is a low-quality,
incomplete, or entirely wrong answer.  The on-chain contract has no
way to evaluate semantic content.

This gap is fundamental to any hash-commitment scheme and represents
the primary trust residual in MAEP's design.  Three mitigation
strategies exist, in increasing order of trustlessness:

\begin{enumerate}
  \item \textbf{LLM Auditor (current).}  The \texttt{AuditorAgent}
    reads both the task specification and the result and issues a
    ruling.  This is practical for open-ended tasks (e.g., text
    analysis, summarization) where quality is inherently subjective,
    but the ruling is non-deterministic and depends on the auditor's
    LLM backend.

  \item \textbf{Verifiable task specification.}  If the task can be
    expressed as a machine-checkable predicate---e.g., ``sort this
    list'', ``execute this SQL query'', ``pass these unit tests''---the
    requester can verify correctness locally and deterministically
    without any auditor.  The protocol would require the provider to
    attach an execution trace or test-pass certificate alongside the
    result hash.  This approach works for structured, computational
    tasks but is inapplicable to generative or open-ended ones.

  \item \textbf{Zero-knowledge proofs (future work).}  A ZK-SNARK or
    ZK-STARK proof can attest that a computation was performed
    correctly without revealing the computation's internals.  For
    example, a provider could prove on-chain that a neural network
    inference on a given input produced a given output, enabling fully
    trustless result acceptance even for generative tasks.  Current ZK
    proving costs for LLM-scale computations remain prohibitive, but
    are an active research area~\cite{wood2014ethereum}.
\end{enumerate}

\textbf{Auditor subjectivity.}
The \texttt{AuditorAgent}'s ruling is LLM-generated and therefore
non-deterministic.  Different runs may produce different rulings for
the same dispute.  For high-stakes use cases, the auditor should be
replaced by a deterministic oracle or a multi-party consensus
mechanism.  The integrity-vs.-correctness limitation above subsumes
this concern for the general case.

\textbf{Testnet vs.\ mainnet.}
Our Gas cost analysis uses Base Sepolia.  Mainnet Base has higher
transaction throughput but also higher absolute gas prices at peak
demand.  The relative Gas efficiency of MAEP's design holds regardless
of absolute price, but practitioners should monitor L2 pricing trends.

\textbf{Limited capability matching.}
The current DISCOVER stage performs exact string matching on
\texttt{task\_types}.  Real deployment would benefit from semantic
matching (embedding similarity) or structured capability schemas.

\subsection{Future Work}

\begin{itemize}
\item \textbf{Forced dispute on silence.}  Currently a requester can
  silently withhold payment until the deadline without triggering
  arbitration.  A future version could require the requester to
  explicitly accept or reject within a response window; silence after
  the window would auto-trigger the \texttt{AuditorAgent} or
  auto-release payment to the provider, closing the liveness gap
  identified in \S\ref{sec:protocol}.

\item \textbf{ZK-proof result verification.} Replace the hash
  commitment with a zero-knowledge proof that the result satisfies a
  verifiable specification (e.g., a SQL query result, a sorting
  correctness predicate), enabling fully trustless result acceptance
  without any auditor.

\item \textbf{Multi-agent task decomposition.} Extend the single
  Requester--Provider bilateral protocol to support a directed acyclic
  graph (DAG) of sub-tasks, where intermediate agents act as both
  providers (upstream) and requesters (downstream).

\item \textbf{Reputation staking and slashing.} Allow third parties to
  stake on an agent's reputation, creating financial incentives for
  honest behavior and a market for reputation insurance.

\item \textbf{Cross-chain deployment.} Abstract the contract layer to
  support deployment on multiple EVM chains simultaneously, with
  cross-chain payment via bridge protocols.
\end{itemize}

%% ============================================================
\section{Conclusion}
\label{sec:conclusion}

We presented MAEP, a Minimum Viable Agent-to-Agent Economic Protocol
that equips autonomous AI agents with on-chain decentralized identity,
trustless escrow payment, and cryptographic result attestation.  By
implementing these primitives as three composable Ethereum smart
contracts and an LLM-agnostic Python SDK, MAEP demonstrates that
trust-minimized agent economies are achievable today without
specialized infrastructure.

Our experimental evaluation shows that MAEP reduces the trust
assumption count to 1 (Ethereum consensus) versus 2--3 for existing
protocols; incurs a total on-chain Gas cost of approximately 408,870
units per task delegation---equivalent to less than one US cent at
current Layer-2 prices; and adds only microsecond-level protocol
overhead above plain function calls.  An end-to-end three-agent
demonstration, including a disputed task that is correctly adjudicated
by an LLM-powered auditor, completes in under five seconds.

We hope MAEP serves as a useful reference point for the emerging field
of agent economics, and invite the community to extend and critique the
design.

%% ============================================================
\section*{Acknowledgments}
The authors thank the open-source communities behind Hardhat, web3.py,
and Ollama for making rapid blockchain prototyping accessible.

%% ============================================================
\begin{thebibliography}{99}

\bibitem{park2023generative}
J.~S. Park, J.~C. O'Brien, C.~J. Cai, M.~R. Morris, P.~Liang, and
M.~S. Bernstein,
``Generative agents: Interactive simulacra of human behavior,''
in \textit{Proc. UIST}, 2023.

\bibitem{w3c_did}
M.~Sporny, D.~Longley, M.~Sabadello, D.~Reed, O.~Steele, and
C.~Allen,
``Decentralized identifiers (DIDs) v1.0,'' W3C Recommendation,
Jul.\ 2022. [Online]. Available: \url{https://www.w3.org/TR/did-core/}

\bibitem{google_a2a}
Google LLC,
``Agent2Agent (A2A) Protocol,'' 2025.
[Online]. Available: \url{https://google.github.io/A2A/}

\bibitem{x402}
Coinbase,
``x402: HTTP Payment Protocol,'' 2025.
[Online]. Available: \url{https://www.x402.org/}

\bibitem{fetchai}
Fetch.ai Foundation,
``Autonomous Economic Agents and the Fetch.ai Network,'' 2024.
[Online]. Available: \url{https://fetch.ai/}

\bibitem{autonolas}
Valory AG,
``Autonolas: On-chain Protocol for Co-owned AI,'' 2024.
[Online]. Available: \url{https://olas.network/}

\bibitem{andres2025}
M.~Andres et al.,
``Towards Multi-Agent Economies: Enhancing the A2A Protocol with
Blockchain-Based Identities and Payments,''
arXiv preprint arXiv:2507.19550, 2025.

\bibitem{consensys_patterns}
ConsenSys,
``Ethereum Smart Contract Security Best Practices: Checks-Effects-Interactions,''
2024. [Online]. Available:
\url{https://consensys.github.io/smart-contract-best-practices/}

\bibitem{wood2014ethereum}
G.~Wood,
``Ethereum: A secure decentralised generalised transaction ledger,''
\textit{Ethereum Project Yellow Paper}, vol.~151, 2014.

\bibitem{nakamoto2008bitcoin}
S.~Nakamoto,
``Bitcoin: A peer-to-peer electronic cash system,'' 2008.
[Online]. Available: \url{https://bitcoin.org/bitcoin.pdf}

\bibitem{szabo1997smart}
N.~Szabo,
``Formalizing and securing relationships on public networks,''
\textit{First Monday}, vol.~2, no.~9, 1997.

\bibitem{buterin2013ethereum}
V.~Buterin,
``Ethereum white paper: A next generation smart contract \& decentralized
application platform,'' 2013.

\end{thebibliography}

\end{document}
